\documentclass[a4paper,12pt]{article}
\usepackage[T2A]{fontenc}
\usepackage[utf8]{inputenc}
\usepackage[english,russian]{babel}
\usepackage{graphicx}
\usepackage{float}
\usepackage{caption}
\usepackage{subcaption}

\title{Лабораторная работа №4 \\ Вставка графики в LaTeX}
\author{Сунь Маосин}
\date{2026}

\begin{document}

\maketitle

\section{Цель работы}

Изучение возможностей LaTeX по вставке и форматированию графических изображений.

\section{Основные команды для вставки графики}

Для вставки графики используется пакет \texttt{graphicx} и команда \texttt{\textbackslash includegraphics}.

\subsection{Простая вставка изображения}

\begin{figure}[H]
  \centering
  \includegraphics{example-image}
  \caption{Простая вставка изображения без изменения размера}
  \label{fig:simple}
\end{figure}

\subsection{Изменение размера изображения}

Можно изменять размер с помощью опций \texttt{width}, \texttt{height} или \texttt{scale}.

\begin{figure}[H]
  \centering
  \includegraphics[width=0.3\textwidth]{example-image}
  \caption{Изображение шириной 0.3 от ширины текста}
  \label{fig:width}
\end{figure}

\begin{figure}[H]
  \centering
  \includegraphics[height=3cm]{example-image}
  \caption{Изображение высотой 3 см}
  \label{fig:height}
\end{figure}

\begin{figure}[H]
  \centering
  \includegraphics[scale=0.5]{example-image}
  \caption{Изображение уменьшено в 2 раза (scale=0.5)}
  \label{fig:scale}
\end{figure}

\subsection{Поворот изображения}

Можно поворачивать изображение с помощью опции \texttt{angle}.

\begin{figure}[H]
  \centering
  \includegraphics[width=0.3\textwidth, angle=45]{example-image}
  \caption{Изображение, повёрнутое на 45 градусов}
  \label{fig:rotate}
\end{figure}

\section{Плавающие окружения (figure)}

\begin{figure}[htbp]
  \centering
  \includegraphics[width=0.5\textwidth]{example-image-a}
  \caption{Пример плавающего окружения figure}
  \label{fig:float}
\end{figure}

\section{Несколько изображений рядом}

\subsection{С помощью minipage}

\begin{figure}[H]
  \centering
  \begin{minipage}{0.3\textwidth}
    \centering
    \includegraphics[width=\textwidth]{example-image}
    \caption{Первое изображение}
    \label{fig:minipage1}
  \end{minipage}
  \hfill
  \begin{minipage}{0.3\textwidth}
    \centering
    \includegraphics[width=\textwidth]{example-image-a}
    \caption{Второе изображение}
    \label{fig:minipage2}
  \end{minipage}
  \hfill
  \begin{minipage}{0.3\textwidth}
    \centering
    \includegraphics[width=\textwidth]{example-image-b}
    \caption{Третье изображение}
    \label{fig:minipage3}
  \end{minipage}
  \caption{Три изображения рядом с помощью minipage}
  \label{fig:minipage}
\end{figure}

\subsection{С помощью subfigure}

\begin{figure}[H]
  \centering
  \begin{subfigure}{0.3\textwidth}
    \centering
    \includegraphics[width=\textwidth]{example-image}
    \caption{Подпись 1}
    \label{fig:sub1}
  \end{subfigure}
  \hfill
  \begin{subfigure}{0.3\textwidth}
    \centering
    \includegraphics[width=\textwidth]{example-image-a}
    \caption{Подпись 2}
    \label{fig:sub2}
  \end{subfigure}
  \hfill
  \begin{subfigure}{0.3\textwidth}
    \centering
    \includegraphics[width=\textwidth]{example-image-b}
    \caption{Подпись 3}
    \label{fig:sub3}
  \end{subfigure}
  \caption{Три изображения с помощью subfigure}
  \label{fig:subfigure}
\end{figure}

\section{Перекрёстные ссылки}

Например, на рисунке~\ref{fig:simple} показана простая вставка изображения.
Рисунок~\ref{fig:width} демонстрирует изменение ширины,
а рисунок~\ref{fig:rotate} — поворот изображения.

На рисунке~\ref{fig:minipage} показаны три изображения рядом,
а на рисунке~\ref{fig:subfigure} — с помощью \texttt{subfigure}.

\section{Вывод}

В ходе работы были изучены:
\begin{itemize}
  \item вставка изображений с помощью пакета \texttt{graphicx};
  \item изменение размера и поворот изображений;
  \item использование плавающего окружения \texttt{figure};
  \item размещение нескольких изображений рядом;
  \item перекрёстные ссылки на рисунки.
\end{itemize}

\end{document}